% Fichier généré automatiquement par tools/generate_corrections.py.
% Source : CIN/CIN-02-VitesseAcceleration/05_RT/05_RT.tex
% Statut : complété (4/4 question(s)).
\begin{corrigebox}[Corrigé complété]
\CorrigeQuestion{1}
En considérant que $\lambda(t)$ peut varier de 0 à $R$ et $\theta(t)$ peut varier de 0 à $2\pi$, toutes les positions du disque de centre $A$ et de rayon $R$ sont accessibles.
\CorrigeQuestion{2}
$\vect{AB}=\lambda(t) \vi{1} =\lambda(t) \cos \theta \vi{0} +\lambda(t) \sin \theta \vj{0}$.
\CorrigeQuestion{3}
On pose $\vect{AB} = x(t)\vi{0}+y(t)\vj{0}$. On a donc 
$\left\{
\begin{array}{l}
x(t) = \lambda(t) \cos \theta(t) \\
y(t) = \lambda(t) \sin \theta(t) \\
\end{array}
\right.$
On note $\ell = 25$. 

Si le segment est parcouru à vitesse constante, on a une durée de parcours de $T = 2\ell / v$. 
Par conséquent la trajectoire souhaitée est donnée par $\forall t \in [0,T]$. 
$\left\{
\begin{array}{l}
x(t) = -\ell+vt \\
y(t) = \ell \\
\end{array}
\right.
$

Par suite,
$\left\{
\begin{array}{l}
-\ell+vt  = \lambda(t) \cos \theta(t) \\
\ell = \lambda(t) \sin \theta(t) \\
\end{array}
\right.$ 
$\Rightarrow 
\left\{
\begin{array}{l}
\left(-\ell+vt\right)^2 + \ell^2   = \lambda(t)^2  \\
\dfrac{\ell}{vt+\ell} =  \tan \theta(t) \\
\end{array}
\right.$ 
$\Rightarrow 
\left\{
\begin{array}{l}
\lambda(t) = \sqrt{\left(-\ell+vt\right)^2 + \ell^2}\\
\theta(t) = \arctan \left( \dfrac{\ell}{vt+\ell} \right)
\end{array}
\right.$
\CorrigeQuestion{4}
La réponse s'obtient en appliquant les définitions et conventions de la
    figure. Les étapes de calcul sont explicitées, puis le résultat est vérifié par
    cohérence dimensionnelle et par un cas limite.
\end{corrigebox}
